\subsection{问题背景}

干燥是决定中药材成品品质的关键工序之一。热风烘干通常包含预热平衡与恒温干燥两个阶段：通过调控烘房的温度和湿度，使圆柱形药材内部的热量向内传递、水分向外迁移，最终把各处干基含水率降到工艺要求以下。工艺参数选择不当会导致干燥效率低、能耗高和成品品质不稳定；完全依赖试验优化则成本高、周期长。因此需要建立能够给出药材内部温度与水分浓度时空分布的数学模型，用数值仿真得到干燥规律。

本题研究对象为长度 $L=25\,\mathrm{cm}$、半径 $R=2\,\mathrm{cm}$ 的近似圆柱药材。烘干开始时药材温度均匀为 $28^\circ\mathrm{C}$，干基含水率均匀为 $2.55\,\mathrm{kg/kg}$。烘房温度与水分浓度随时间的变化由附件给出；第四问另给出烘干过程中药材半径随时间的变化。相关热物性、扩散系数及对流换热、对流传质系数由各问附录给出。

\subsection{问题内容}

\textbf{问题一：}在预热平衡阶段，采用附录给出的常物性参数，建立药材温度 $T(r,t)$ 与干基含水率 $C(r,t)$ 的变化规律。按题目表格给出 $100,300,600,900,1200,1500,1800\,\mathrm{s}$ 时、到中心距离 $0,0.5,1,1.5,2\,\mathrm{cm}$ 处的温度和含水率，并保存 $1800\,\mathrm{s}$ 内每隔 $1\,\mathrm{s}$、半径每隔 $0.1\,\mathrm{cm}$ 的完整结果。

\textbf{问题二：}烘干过程一般持续 $2$--$3$ 天，预热与恒温阶段的物性不同。为简化问题，相关经验公式统一采用另一组随含水率和温度变化的公式，建立整个烘干过程前 $3\,\mathrm{h}$ 的温度与含水率模型。按表格给出每隔 $0.5\,\mathrm{h}$、上述五个半径处的结果，并保存每隔 $1\,\mathrm{s}$、半径每隔 $0.1\,\mathrm{cm}$ 的完整结果。

\textbf{问题三：}按烘干要求，药材各处水分浓度应低于 $0.15\,\mathrm{kg/kg}$。确定烘干所需时间，并按表格给出每隔 $6\,\mathrm{h}$ 及结束时刻、五个半径处的含水率。

\textbf{问题四：}实际烘干中药材因失水发生尺寸变化。根据给定的半径--时间数据，更换相应经验物性后重新确定烘干时长，并按动态半径输出含水率表格。
